\documentclass[12pt,a4paper]{article}
\usepackage[utf8]{inputenc}
\usepackage[spanish,es-nodecimaldot,es-tabla]{babel}
\usepackage{graphicx}
\usepackage{float}
\usepackage{geometry}
\usepackage{hyperref}
\usepackage{fancyhdr}
\usepackage{titlesec}
\usepackage{booktabs}
\usepackage{xcolor}
\usepackage{verbatim}
\usepackage{mathptmx}
\usepackage{setspace}
\usepackage{xurl}
\usepackage{amsmath}

\geometry{top=2.5cm, bottom=2.5cm, left=3cm, right=3cm}
\hypersetup{colorlinks=true, linkcolor=black, urlcolor=blue, citecolor=black}
\graphicspath{{../resultados/figuras/}{figuras_codigo/}}

\pagestyle{fancy}
\fancyhf{}
\fancyhead[L]{\small Colonia de hormigas aplicada al TSP}
\fancyhead[R]{\thepage}
\renewcommand{\headrulewidth}{0.4pt}
\setlength{\headheight}{15pt}

\titleformat{\section}{\normalfont\Large\bfseries}{\thesection}{1em}{}
\titleformat{\subsection}{\normalfont\large\bfseries}{\thesubsection}{1em}{}
\titleformat{\subsubsection}{\normalfont\normalsize\bfseries}{\thesubsubsection}{1em}{}

\begin{document}

\begin{titlepage}
\thispagestyle{empty}
\centering

\vspace*{0.5cm}
\includegraphics[width=0.30\textwidth]{sergio_logo.png}\\[0.4cm]

{\fontsize{13}{15}\selectfont\textbf{UNIVERSIDAD SERGIO ARBOLEDA}}\\[2.5cm]

{\fontsize{14}{17}\selectfont\textbf{COLONIA DE HORMIGAS APLICADA AL PROBLEMA\\[0.2cm]
DEL VIAJANTE (TSP): INSTANCIAS DE 20, 2 000 Y 200 000 CIUDADES}}\\[3.5cm]

{\fontsize{12}{14}\selectfont\textbf{Inteligencia Artificial}}\\[1.8cm]

{\fontsize{12}{14}\selectfont Mario Jiménez}\\[0.3cm]
{\fontsize{12}{14}\selectfont Juan David Andrade}\\[0.3cm]
{\fontsize{12}{14}\selectfont Camilo Andrés Díaz García}\\[0.3cm]

\vfill

\begin{center}
    \textit{Universidad Sergio Arboleda}\\
    \textit{Facultad de Sistemas e Ingeniería}\\
    \textit{Bogotá D.C., 2026}
\end{center}

\end{titlepage}

\newpage
\tableofcontents
\newpage

\onehalfspacing

% ====================================================================
\section{Introducción}
\emph{Pendiente: se completa en la Fase 5, con el contexto del taller, el
equipo y el objetivo de ahorrar tiempo y memoria frente al espacio de
búsqueda factorial del TSP.}

% ====================================================================
\section{Marco teórico}
\emph{Pendiente: estigmergia, regla de transición, evaporación, depósito,
exploración y explotación, con las fórmulas de las diapositivas 14--19 (ver
`docs/bitacora.md` y el código fuente para las expresiones ya validadas).}

% ====================================================================
\section{Descripción del problema y espacio factorial}
\emph{Pendiente: TSP simétrico con origen fijo, $(n-1)!/2$ recorridos
posibles, ejemplos numéricos para $n=20$, $2\,000$ y $200\,000$ vía
\texttt{lgamma} (Fase 4).}

% ====================================================================
\section{Diseño de la solución}

\subsection{Elección de lenguaje}
El proyecto se implementa en C++20 con OpenMP; no se implementó nada en
Python, C ni Rust, y no se hicieron comparaciones de lenguaje. La
justificación es breve y teórica, no empírica:

Python se descarta por tres razones estructurales: la sobrecarga del bucle
de evaluación de bytecode del intérprete de CPython en cada operación
aritmética del paso de una hormiga; el acceso irregular a memoria, que
Python (incluso con \texttt{numpy} en presencia de lógica de control como la
ruleta) no permite controlar con la precisión de un arreglo contiguo de
$n \cdot K$ elementos; y la falta de paralelismo real dentro de un proceso,
por el \emph{Global Interpreter Lock} (GIL) de CPython, que impide que dos
hilos ejecuten bytecode de Python al mismo tiempo, justo cuando la parte
costosa del algoritmo (la construcción de tours por hormigas independientes)
es la que se querría paralelizar con hilos.

C++ ofrece en cambio control de memoria (arreglos contiguos, sin recolector
de basura), OpenMP como mecanismo directo de paralelismo por hilos, y un
código base ya validado (\texttt{aco\_tsp.cpp}, corrido en Linux con g++ 13
antes de este repositorio). C y Rust no se descartan por incapacidad
técnica del enunciado, sino porque no aportaban nada sobre partir de ese
código ya validado en C++. Ver \texttt{docs/decisiones\_diseno.md} para la
tabla completa de la decisión.

\subsection{Número de agentes}
\label{sec:numero-agentes}
Se adopta $m = \min(n, 20\,000)$: $n=20$ con $m=20$; $n=2\,000$ con
$m=2\,000$; $n=200\,000$ con $m=20\,000$ (las $200\,000$ ciudades se
ejecutan siempre completas).

\paragraph{Convención de la literatura.} Dorigo y Stützle (2004, \emph{Ant
Colony Optimization}, MIT Press) documentan, en la tabla de parámetros
recomendados para AS y sus variantes sin búsqueda local (p.~71, Box~3.1),
que el número de hormigas se fija en $m=n$; en las pp.~103--104 explican que
esto hace el requisito de memoria $O(n^2)$ (${\approx}32n^2$ bytes); y en la
p.~112 atribuyen la convención $m \propto n$ a los artículos originales de
Dorigo, Maniezzo y Colorni (1991, 1996). No se cita una página de esos
artículos originales por no haberlos verificado directamente.

\paragraph{Costo y extrapolación.} Con lista de candidatos de tamaño $K$,
una iteración cuesta $O(m n K)$; con $m=n$ pasa a $O(n^2 K)$. El tiempo
medido en esta máquina para una iteración en $n=200\,000$, $K=8$
(\texttt{resultados/costo\_m\_n200000.csv}) fue de 1,099~s ($m=200$),
13,056~s ($m=2\,000$) y 144,850~s ($m=20\,000$). El ajuste lineal por
mínimos cuadrados da $t(m) \approx 0{,}007285\,m - 0{,}909$ segundos, con
$R^2 = 0{,}999947$. \textbf{Como estimación por extrapolación lineal, no
como corrida real}, una iteración con $m=n=200\,000$ tomaría
${\approx}1\,456$ s (${\approx}24{,}3$ min), cifra consistente con la
estimación independiente de la Fase~1 (${\approx}24{,}5$ min).

\paragraph{Por qué el tope de 20 000.} Es una decisión de presupuesto de
cómputo de esta máquina (Fase~0), no una propiedad del algoritmo: diez
iteraciones con $m=n=200\,000$ tomarían más de 4 horas (extrapolado),
mientras que con $m=20\,000$ toman del orden de 25 minutos por semilla,
medido directamente. Además, la feromona se actualiza una sola vez por
iteración sin importar cuántas hormigas corrieron, así que agregar hormigas
no agrega actualizaciones de feromona, solo muestras adicionales antes de
esa única actualización.

\paragraph{Respaldo empírico.} La calidad contra $m$ en $n=2\,000$ (50
iteraciones fijas, 5 semillas, \texttt{resultados/barrido\_m.csv}) muestra
rendimientos decrecientes marcados: de $m=10$ a $m=20\,000$ (2\,000 veces
más hormigas) la longitud media mejora solo 2,4\,\% (39,84 a 38,89), con un
costo de cómputo casi 1\,080 veces mayor. La comparación a igual presupuesto
de cómputo en $n=200\,000$ ($m=2\,048$ contra $m=20\,000$, mismo tiempo de
reloj) queda \textbf{pendiente} de ejecutarse (corrida larga, con
confirmación previa del estudiante); esta sección se actualiza cuando esa
corrida cierre. Con los datos disponibles hasta ahora, la saturación de
calidad con el número de iteraciones fijo sugiere que el tope de 20\,000 no
responde a una necesidad de calidad, sino al presupuesto de tiempo — pero
esa conclusión se confirma o se corrige con la comparación pendiente.

\paragraph{Aclaración.} 20\,000 es una decisión de ingeniería para esta
máquina y este equipo, no un óptimo teórico. Dorigo y Stützle (2004, p.~217,
sección 5.7.6) señalan que "el mejor valor de $m$ es función del algoritmo
ACO particular y de la clase de problemas atacados, y la mayoría de las
veces debe fijarse experimentalmente".

\subsection{Estructuras de datos y estrategias de ahorro}
\emph{Pendiente: lista de $K$ vecinas por ciudad, rejilla espacial,
renumeración por celda, buffers por hilo sin almacenar el tour, pesos
precalculados por iteración (ver \texttt{docs/estrategias\_ahorro.md}).}

% ====================================================================
\section{Implementación}
\emph{Pendiente: portabilidad a Windows, \texttt{--input} TSPLIB,
\texttt{--tiempo\_max}, versión densa de referencia, pruebas automáticas
(Fase 1, ver \texttt{docs/bitacora.md}).}

% ====================================================================
\section{Metodología experimental}
\emph{Pendiente: barrido de parámetros con 5 semillas, criterio de parada,
estudios de escalamiento (Fase 2).}

% ====================================================================
\section{Resultados por instancia}
\emph{Pendiente: tablas de longitud, tiempo y memoria por instancia (20,
2\,000, 200\,000), con media, desviación, mínimo y máximo (Fase 2 y
\texttt{docs/resultados.md}).}

% ====================================================================
\section{Análisis de tiempo y memoria}
\emph{Pendiente: complejidad y comprobación empírica (Fase 4).}

% ====================================================================
\section{Discusión}
\emph{Pendiente.}

% ====================================================================
\section{Limitaciones}
\emph{Pendiente: ACO no garantiza el óptimo; con $n=200\,000$ y pocas
iteraciones el resultado queda cerca del vecino más cercano; la lista de
candidatos restringe el espacio de búsqueda; la instancia aleatoria uniforme
no es una instancia TSPLIB real.}

% ====================================================================
\section{Conclusiones}
\emph{Pendiente.}

% ====================================================================
\section{Referencias}
\emph{Pendiente: Hurbans (2020, cap. 6); Dorigo, Maniezzo y Colorni (1996);
Dorigo y Stützle (2004).}

\end{document}
